\section{Metodologia: SFC-ABM jako uzupełnienie DSGE}

Metoda łączy dwie tradycje: stock-flow consistent modelling oraz agent-based computational economics. Podejście SFC, rozwinięte m.in. przez Godleya i Lavoie \parencite{godley2007} i omawiane szerzej w literaturze postkeynesowskiej \parencite{caverzasi2015}, wymaga, aby każdy przepływ miał drugą stronę księgową, a każdy zasób wynikał z akumulacji przepływów. Z kolei ABM, opisany w tradycji Tesfatsion \parencite{tesfatsion2006}, traktuje gospodarkę jako system heterogenicznych agentów podejmujących decyzje w warunkach ograniczonej racjonalności i wytwarzających dynamikę makroekonomiczną całego układu. Połączenie obu tradycji jest obecne w literaturze SFC-ABM i makroekonomii złożoności, w której nacisk pada na bilanse, niejednorodność oraz endogenne procesy dostosowawcze \parencite{caiani2016,dosi2010,delligatti2011,dafermos2017}.

Połączenie SFC i ABM jest użyteczne tam, gdzie pytanie dotyczy ścieżki przejścia, bilansów i niejednorodności. W takim modelu gospodarstwa, firmy, banki, państwo i sektor zewnętrzny są powiązane przepływami pieniężnymi. Najkrótszym zapisem tej logiki jest równanie bilansu sektorowego, często wiązane z tradycją Godleya:

\[
(S - I) + (T - G) + (M - X) = 0,
\]

gdzie \(S - I\) oznacza saldo sektora prywatnego, \(T - G\) saldo sektora publicznego, a \(M - X\) saldo sektora zagranicznego względem gospodarki krajowej. Nie jest to założenie behawioralne, lecz tożsamość rachunkowa: nadwyżka jednego sektora musi mieć odpowiednik w deficycie innego sektora. Jeżeli państwo zwiększa transfery, musi pojawić się odpowiedni zapis po stronie dochodów gospodarstw, deficytu, długu, podatków albo sektora zewnętrznego. Takie ograniczenie nie jest detalem technicznym, lecz warunkiem poprawności eksperymentu.

SFC-ABM nie zastępuje DSGE. DSGE lepiej nadaje się do formalnych porównań dobrobytowych i analizy optymalnych polityk w zdefiniowanych strukturach preferencji oraz oczekiwań. SFC-ABM lepiej nadaje się do analizy płynności, zadłużenia, bankructw, ograniczeń bankowych, kursu walutowego i mikroheterogeniczności firm. Dlatego w tej pracy oba podejścia są używane łącznie.

Ważna jest także falsyfikowalność. Eksperyment obliczeniowy nie powinien być wyłącznie opisem słownym. Powinien być możliwy do powtórzenia przez zewnętrznego czytelnika. Z tego powodu narzędzie użyte w pracy, \amor{}, jest otwartym repozytorium, a scenariusz BDP został zapisany jako jawny plik zmian i odtworzony w aneksie replikacyjnym.
